\chapter{End-to-End Software and MoveApps Workflow}
\label{chap:software}

\section{Workflow as the central contribution}

\begin{itemize}
  \item The software connects research prototypes and theoretical algorithms to an executable workflow for real movement data.
  \item Every component can be used independently, replaced by an alternative, or composed into the complete chain.
  \item The full route through the system is:
  \begin{itemize}
    \item validate and annotate observations;
    \item project trajectory and environment into compatible metric grids;
    \item split gaps and annotate movement structure;
    \item fit state-specific movement kernels;
    \item estimate or specify terrain effects;
    \item construct endpoint-conditioned paths;
    \item calculate UDs and visit probabilities;
    \item export machine-readable artifacts and optional Leaflet visualizations.
  \end{itemize}
\end{itemize}

\section{Components and responsibilities}

\begin{itemize}
  \item \textbf{Environmental Land Cover Annotation} (\code{cma\_landcover}/\code{environmentcma}): LOCAL WorldCover or LONG\_RANGE land/water mode, per-animal projection, terrain fields, grids, and optional GeoTIFF archive.
  \item \textbf{Behavioral Segmentation} (\code{hmm\_moveapp}/\code{hmmcma}/\code{segmentationcma}): criteria rules or data-driven states with common annotation and diagnostic fields.
  \item \textbf{Kernel Construction} (\code{kernelcma}): Brownian/correlated state kernels, density estimator configuration, physical time interval, automatic probability-mass support, and serialized kernel metadata.
  \item \textbf{Random Walks MoveApp}: endpoint-conditioned, state-dependent, and terrain-aware path generation with user-facing animal type, barrier mode, spatial resolution, and movement-model settings.
  \item \textbf{Utilization Distribution MoveApp}: reuse existing states, calculate conditioned UDs, optionally sample paths, and return spatial artifacts.
  \item \textbf{Research integration package} (\code{/home/omar/PycharmProjects/RW-Python-gitlab/}): shared Python interfaces, C/CUDA bindings, experiments, reference comparisons, and regression tests.
\end{itemize}

\section{Cross-component contract}

\begin{itemize}
  \item Preserve animal identifier, timestamp, geometry, source CRS, projected coordinates, grid coordinates, terrain, state, and original columns.
  \item Attach metadata for units, projection, raster product, grid origin/cell size, class/state mappings, kernel interval/support, random seeds, component version, and warnings.
  \item Reject missing or incompatible prerequisites with an actionable error.
  \item Avoid hidden inference:
  \begin{itemize}
    \item the UD app accepts an explicit state column rather than silently refitting behavior;
    \item local and long-range environmental modes communicate different assumptions about resolution and coverage;
    \item barrier choices are explicit and taxon-dependent;
    \item sampled paths remain distinguishable from observations.
  \end{itemize}
\end{itemize}

\section{Python/native boundary}

\begin{itemize}
  \item Python owns ingestion, geospatial preprocessing, model configuration, orchestration, diagnostics, and visualization.
  \item C/OpenMP/CUDA own performance-critical transition processing and dynamic-programming operations.
  \item Bindings validate array shape, dtype, memory layout, dimensions, model parameters, kernel mappings, and ownership/lifetime.
  \item Errors must propagate across the boundary without returning partially valid trajectories.
  \item Deterministic seeds and stable serialization support regression testing and reproducibility.
\end{itemize}

\section{User-facing design}

\begin{itemize}
  \item Expose choices users can interpret ecologically: study range, segmentation method, features/rules, animal type, barrier mode, kernel family, time factor where valid, resolution, and number of sampled paths.
  \item Derive internal quantities such as exact grid dimensions, kernel cell radius, and minimum feasible transition count where possible.
  \item Return concise diagnostics explaining automatic choices and adaptations.
  \item Package intermediate artifacts only when they support inspection, reuse, or reproducibility.
  \item Include one complete MoveApps workflow screenshot/configuration and one step-by-step artifact trace in the final version.
\end{itemize}

\section{Testing and provenance}

\begin{itemize}
  \item Unit tests: coordinates, class mappings, features, normalization, path rasterization, parameter validation, and serialization.
  \item Property tests: probability conservation, fixed endpoints, barrier compliance, and equivalence of representations.
  \item Integration tests: each adjacent MoveApps contract and the full workflow.
  \item Regression tests: representative terrestrial, marine, airborne, local, and long-range studies.
  \item Numerical tests: serial versus OpenMP/CUDA outputs within declared tolerances.
  \item Provenance: immutable repository revisions, package versions, app settings, data identifiers, environmental asset identifiers, seeds, and raw output paths.
\end{itemize}
